\section{Introduction}
\label{sec:introduction}

Let $[N]=\{1,\ldots,N\}$, and let $\mu_p$ be the probability measure
on $2^{[N]}$ under which each coordinate is present independently with
probability $p\in(0,1)$. For a family $\mathcal D\subseteq2^{[N]}$ and
an integer $k\ge1$, define
\begin{equation}
\label{eq:exceptional-definition}
E_k(\mathcal D)
=\bigl\{S\subseteq[N]:
  S\nsubseteq D_1\cup\cdots\cup D_k
  \text{ for all }D_1,\ldots,D_k\in\mathcal D\bigr\}.
\end{equation}
The members in a tuple may repeat and need not be disjoint.

A family $\mathcal G\subseteq2^{[N]}$ \emph{covers} a family
$\mathcal A$ if every $S\in\mathcal A$ contains some $I\in\mathcal G$.
For $\rho\in(0,1)$, write
\[
c_\rho(\mathcal G)=\sum_{I\in\mathcal G}\rho^{|I|}.
\]
We say that $\mathcal A$ is \emph{$p$-small} if it has a cover
$\mathcal G$ with $c_p(\mathcal G)\le1/2$.
Since $\mu_p\{S:I\subseteq S\}=p^{|I|}$, such a cover witnesses
$\mu_p(\mathcal A)\le1/2$ by a union bound.

The name reflects the analogy with Talagrand's Gaussian convexity
question~\cite[Conjecture~2.1]{Talagrand10}: can a bounded number of
sums of a compact balanced set of sufficiently large Gaussian measure
contain a convex set of Gaussian measure at least $1/2$, uniformly in
the dimension? Here balanced means closed under multiplication by
scalars of absolute value at most one. In the discrete problem, unions
replace sums, and containment covers describe the exceptional sets.
Talagrand formulated the discrete-convexity conjecture in
\cite[Conjecture~7.1]{Talagrand10}: a fixed number of members of a
sufficiently likely family should cover every set outside a $p$-small
exceptional family. The same original-density formulation appears as
Research Problem~13.3.2 in~\cite{Talagrand21}, whose
Proposition~13.3.3 explains its connection to positive selector processes.
We prove this conjecture.

\begin{theorem}
\label{thm:main}
Set $k=2^{75}$. For every integer $N\ge1$, every $p\in(0,1)$, and every
family $\mathcal D\subseteq2^{[N]}$,
\[
\mu_p(\mathcal D)\ge1-\frac1k
\quad\Longrightarrow\quad
E_k(\mathcal D)\text{ is }p\text{-small}.
\]
\end{theorem}

The constant is independent of $N$, $p$, and $\mathcal D$.
The family $\mathcal D$ is arbitrary; no monotonicity is assumed.

\paragraph{Positive selectors.}
The connection to positive selector processes
\cite[Section~7]{Talagrand10}
gives a concrete consequence of the theorem, recovering the
positive-selector conclusion proved by Park and Pham
\cite[Theorem~1.2]{ParkPham24}. For a nonempty family
$T\subseteq[0,\infty)^N$, define
\[
\phi(S)=\sup_{t\in T}\sum_{i\in S}t_i,
\qquad M=\mathbb E_{\mu_p}\phi,
\]
and suppose $0<M<\infty$. Then
$\{S:\phi(S)\ge k^2M\}$ is $p$-small for $k=2^{75}$.
Indeed, Markov's inequality gives
$\mu_p\{D:\phi(D)<kM\}\ge1-1/k$.
Nonnegativity of the coordinates makes $\phi$ monotone and subadditive
under unions, so a set with $\phi(S)\ge k^2M$ cannot be contained in
the union of $k$ such $D$'s. Theorem~\ref{thm:main} supplies its cover.
Thus the large-value event has explicit containment witnesses, beyond
the small probability already supplied by Markov's inequality.

\paragraph{Related developments.}
Hua, Song, and Tudose resolve Talagrand's Gaussian convexity problem and
derive a discrete conclusion with three unions and covering density
$p^C$, for a universal constant $C$, when the starting family is
sufficiently likely~\cite[Corollaries~1.3--1.4]{HuaSongTudose26}.
Ascoli, He, Park, and Talagrand give a threshold reformulation of discrete
convexity and prove cases for several graph-containment
properties~\cite[Proposition~1.5 and Theorems~1.10--1.13]{AscoliEtAl26}.

There is also a fractional version of smallness: instead of selecting
generators, one assigns nonnegative coefficients $a_I$ satisfying
\[
\sum_{I\subseteq S}a_I\ge1\quad(S\in\mathcal A),
\qquad
\sum_{I\subseteq[N]}p^{|I|}a_I\le\frac12.
\]
Talagrand introduced this relaxation alongside integral
smallness~\cite[Sections~6 and~7]{Talagrand10}.
Li resolves the fractional discrete-convexity problem through a quantitative
obstruction to probability measures whose containment probabilities
are bounded by $p^{|I|}$~\cite[Theorem~1.9]{Li26}.
Combining fractional information with rounding, Park obtains an integral
cover of $E_2(\mathcal D)$ at density
\[
\frac{p}{C\max\{1,\log\log(1/p)\}}
\]
whenever $\mu_p(\mathcal D)>1/2$, for a universal
$C$~\cite[Corollary~3.5]{Park26}.
Combining Li's result with the companion constant-factor rounding theorem
gives the following two-union tradeoff.

\begin{corollary}[Two unions at constant density loss]
\label{cor:two-union}
There is a universal constant $C\ge1$ such that, for every integer $N\ge1$,
every $p\in(0,1)$, and every family $\mathcal D\subseteq2^{[N]}$,
\[
\mu_p(\mathcal D)>\frac12
\quad\Longrightarrow\quad
E_2(\mathcal D)\text{ is }(p/C)\text{-small}.
\]
\end{corollary}

\begin{proof}
If $E_2(\mathcal D)=\varnothing$, the empty cover suffices. Otherwise
it is a nonempty proper increasing family: $\mathcal D\ne\varnothing$,
so $\varnothing\notin E_2(\mathcal D)$, and the defining obstruction
is preserved by taking supersets. By Li's result
\cite[Corollary~3.1]{Li26} in Park's fractional-cover formulation
\cite[Theorem~3.4]{Park26}, this family has a fractional cover of cost
at most $1/2$ at $r=p/2$. Capping each coefficient at one preserves
coverage and does not increase cost. The companion expectation-threshold
theorem~\cite[Theorem~1.1]{OpenAIThreshold26} then gives an integral cover
of cost at most $1/2$ at
$r/(25\cdot512^4)$. Thus $C=50\cdot512^4$ suffices.
\end{proof}

Theorem~\ref{thm:main} retains the original cost parameter $p$ while
allowing a larger universal number of unions.

\paragraph{Proof strategy and method.}
The main task is to turn the obstruction to covering a set by unions
into an inexpensive family of integral generators. We first work at a
fixed fraction of the density. An orthogonal expansion gives nonnegative
weights on subsets with a global bound on their weighted sum. A signed
product on an array of 32 rows turns the failure of a union cover into
a lower bound on the sixteenth powers of those same weights. Its column
factor vanishes when all rows miss that coordinate, so the identity
follows from pointwise vanishing. The two-row version is Li's signed
kernel after multiplication by the Bernoulli masses~\cite[Section~2]{Li26}.
This method has antecedents
in the signed product operators and biased Fourier analysis of
Friedgut~\cite[Section~2.2]{Friedgut08} and Ellis, Filmus, and
Friedgut~\cite[Sections~2.2 and~4.1]{EllisFilmusFriedgut12}.

Large individual weights give cheap generators directly. If a set
contains none of these generators, the higher-power bound forces it
to contain many subsets of one size and comparable weight. An elementary
covering lemma handles each such class: it selects subsets of large
weighted degree, independently samples unions of pairs of the remaining
edges, and adds the residual configurations. Weighted-degree control
bounds the cost of extending a tuple, while a long tuple must avoid
quadratically many independent pair choices to survive. Random unions
also enter the covering constructions of Fischer and
Person~\cite[Sections~1.1 and~2]{FischerPerson25}; the lemma here is proved
in full.

Section~\ref{sec:covering} establishes this covering lemma, and
Section~\ref{sec:weights} uses it with the signed identity to cover
$E_{32}(\mathcal F)$ at density $q/2^{70}$ whenever
$\mu_q(\mathcal F)\ge1/2$. The final step, in
Section~\ref{sec:density}, recovers the original density. We couple
dependent rows so that each entire row has the original Bernoulli product
law and their union has an exactly prescribed larger density. The
probability that all rows belong to the starting family is controlled by
a union bound. This exact coupling is what permits arbitrary families
and returns the covering cost to $p$.
